\section{Formal Protocol and Statistical Estimand}
\label{app:formal}

This appendix gives the complete probability model and proofs underlying the
SAFE-ALPHA certificate.  Time is discrete.  Let
\((\Omega,\mathcal F,(\mathcal F_t)_{t\geq 0},P)\) be a filtered probability
space, where \(\mathcal F_t\) is the \emph{global} information available to
the research system by time \(t\).  It contains all market observations,
data releases, agent inputs and outputs, candidate specifications, execution
records, evidence streams, and gate decisions observed by then.  For a
stopping time \(\tau\), write
\[
  \mathcal F_\tau
  =\{B\in\mathcal F:B\cap\{\tau\leq t\}\in\mathcal F_t
    \text{ for every }t\}.
\]

Candidate \(j\) arrives at an \((\mathcal F_t)\)-stopping time \(\tau_j\).
Its executable specification \(A_j\) is
\(\mathcal F_{\tau_j}\)-measurable.  The specification includes the signal
formula, universe, data vintage, execution lag, missing-data convention,
position and turnover definitions, cost rule, score transformation, and
evidence rule.  We index only candidates that enter the protocol, so
\(\tau_j<\infty\) almost surely.  There may be infinitely many candidates
over the life of the system, but only finitely many may have arrived by any
finite calendar time.  A substantive revision of \(A_j\) receives a new
index, a new proposal time, and unit initial evidence.

For \(t>\tau_j\), let \(X_{j,t}\) be the score produced by the committed
strategy.  It is \(\mathcal F_t\)-measurable, and for the linear betting
construction below we assume
\begin{equation}
  -1\leq X_{j,t}\leq 1\quad\text{almost surely}.
  \label{eq:app-bounded}
\end{equation}
For fixed \(P\), let \(N_j\in\{0,1\}\) indicate that the randomly selected
candidate is a true null.  We require \(N_j\) to be
\(\mathcal F_{\tau_j}\)-measurable and take the following selective
conditional-drift restriction as the primitive null condition:
\begin{equation}
  N_j\mathbf 1\{\tau_j<t\}
  \mathbb E_P[X_{j,t}\mid\mathcal F_{t-1}]\leq 0
  \quad\text{almost surely for every }t.
  \label{eq:app-selective-null}
\end{equation}
Thus the identity of a null may be selected from the global past.  What is
excluded is selection with information from the outcome that will test it.
One formal construction starts with a countable or standard-Borel strategy
grammar \(\mathcal S\).  Each \(a\in\mathcal S\) has a null model
\(\mathcal P_a\), \(A_j\) is selected at \(\tau_j\), and
\(N_j(P)=\mathbf 1\{P\in\mathcal P_{A_j}\}\).  In that construction,
membership in \(\mathcal P_a\) must imply
Eq.~\eqref{eq:app-selective-null} uniformly over every admissible adaptive
history that can lead to \(A_j=a\).  Marginal mean zero for each strategy
when prespecified does not suffice.

\subsection{What is certified}
\label{app:estimand}

The theorem certifies the estimand represented by \(X_{j,t}\), not an
unstated raw-return estimand.  To make this distinction explicit, suppose
the committed position earns gross excess return \(G_{j,t}\), its turnover is
\[
  U_{j,t}=\frac12\lVert w_{j,t}-w_{j,t^-}\rVert_1,
\]
and the ledger charges \(c_{j,t}\) per unit of turnover.  The tested net
return is
\begin{equation}
  R^{\mathrm{net}}_{j,t}=G_{j,t}-c_{j,t}U_{j,t}.
  \label{eq:app-net-return}
\end{equation}
The position, execution rule, and accounting convention must be fixed before
the corresponding return is realized.  A conservative cost schedule may be
fixed at proposal time; it may also be updated predictably, with
\(c_{j,t}\in\mathcal F_{t-1}\).  Actual fills and realized fees can instead
enter Eq.~\eqref{eq:app-net-return} through a precommitted accounting rule.
Choosing the cost charge after inspecting \(G_{j,t}\) is not admissible.

In the empirical protocol, a predictable scale \(b_{j,t}>0\) is formed from
lagged volatility and
\begin{equation}
  X_{j,t}
  =\operatorname{clip}\!\left(
      \frac{R^{\mathrm{net}}_{j,t}}{b_{j,t}},-1,1\right),
  \qquad
  b_{j,t}=4\max\{\widehat\sigma_{j,t-1},10^{-6}\},
  \label{eq:app-score}
\end{equation}
where \(\widehat\sigma_{j,t-1}\) is estimated from 63 lagged observations.
A lagged expanding estimate is used during initialization.  The null concerns
the conditional mean of this bounded net score.  In general,
\[
  \mathbb E[R^{\mathrm{net}}_{j,t}\mid\mathcal F_{t-1}]\leq 0
  \ \not\Longrightarrow\
  \mathbb E[\operatorname{clip}(R^{\mathrm{net}}_{j,t}/b_{j,t},-1,1)
       \mid\mathcal F_{t-1}]\leq 0.
\]
Consequently, the certificate must not be reported as a test of raw mean
return.  A raw-mean interpretation requires an actually bounded payoff
implemented through position limits, a known one-sided payoff bound, or a
different e-process justified by a stated conditional moment-generating
function or variance assumption.  The multiple-testing theorem below is
agnostic about that choice: it uses only the selective e-validity property
proved next.

\section{Global E-Validity After Adaptive Proposal}
\label{app:evalidity}

Let \(\lambda_{j,t}\in[0,1]\) be
\(\mathcal F_{t-1}\)-measurable.  The bet may use every element of the global
past, including other candidates' histories, but not the current payoff.
Define
\begin{equation}
  E_{j,t}=
  \begin{cases}
    1, & t\leq\tau_j,\\[1mm]
    \displaystyle\prod_{s=\tau_j+1}^{t}
       (1+\lambda_{j,s}X_{j,s}), & t>\tau_j .
  \end{cases}
  \label{eq:app-eprocess}
\end{equation}
Every factor is nonnegative by Eq.~\eqref{eq:app-bounded}.

\begin{lemma}[Global e-validity]
\label{lem:app-selective-e}
Suppose Eq.~\eqref{eq:app-selective-null} holds.  Let \(S\geq\tau_j\) be any
global stopping time, possibly unbounded or infinite, and define its killed
stopping value by
\[
  \overline E_{j,S}
  =E_{j,S}\mathbf 1\{S<\infty\},
\]
with the product in Eq.~\eqref{eq:app-eprocess} evaluated only on
\(\{S<\infty\}\).  Then
\begin{equation}
  \mathbb E_P[N_j\overline E_{j,S}\mid\mathcal F_{\tau_j}]
  \leq N_j\quad\text{almost surely}.
  \label{eq:app-selective-e}
\end{equation}
In particular,
\(\mathbb E_P[N_j\overline E_{j,S}]\leq\mathbb E_P[N_j]\).
\end{lemma}

\begin{proof}
First consider one post-proposal update.  On the event \(N_j=1\) and
\(\tau_j<t\), predictability and
Eq.~\eqref{eq:app-selective-null} give
\begin{align}
 \mathbb E_P[E_{j,t}\mid\mathcal F_{t-1}]
 &=E_{j,t-1}
   \mathbb E_P[1+\lambda_{j,t}X_{j,t}\mid\mathcal F_{t-1}]
 \notag\\
 &=E_{j,t-1}\left\{1+\lambda_{j,t}
   \mathbb E_P[X_{j,t}\mid\mathcal F_{t-1}]\right\}
 \leq E_{j,t-1}.
 \label{eq:app-one-step}
\end{align}
When \(N_j=0\), multiplication by \(N_j\) makes both sides zero.

The random starting time requires a short argument.  Set
\(\mathcal G_{j,n}=\mathcal F_{\tau_j+n}\) for \(n\geq0\).  Because
\(\tau_j+n\) is a stopping time, this is a filtration, and
\(N_jE_{j,\tau_j+n}\) is \(\mathcal G_{j,n}\)-measurable.  If
\(B\in\mathcal G_{j,n-1}\), then
\(B\cap\{\tau_j=r\}\in\mathcal F_{r+n-1}\).  Partitioning over the possible
values of \(\tau_j\) and using Eq.~\eqref{eq:app-one-step} yields
\begin{align*}
 \mathbb E_P[\mathbf 1_BN_jE_{j,\tau_j+n}]
 &=\sum_{r\geq0}
   \mathbb E_P[\mathbf 1_{B\cap\{\tau_j=r\}}
                    N_jE_{j,r+n}]\\
 &\leq\sum_{r\geq0}
   \mathbb E_P[\mathbf 1_{B\cap\{\tau_j=r\}}
                    N_jE_{j,r+n-1}]\\
 &=\mathbb E_P[\mathbf 1_BN_jE_{j,\tau_j+n-1}].
\end{align*}
It follows that
\[
  (N_jE_{j,\tau_j+n})_{n\geq0}
\]
is a nonnegative \((\mathcal G_{j,n})\)-supermartingale.  Its initial value
is \(N_j\).

The shifted time \(R=S-\tau_j\) is a nonnegative
\((\mathcal G_{j,n})\)-stopping time: for every \(n\),
\(\{R\leq n\}=\{S\leq\tau_j+n\}\in\mathcal G_{j,n}\).  Conditional optional
sampling at \(R\wedge m\) therefore gives
\begin{equation}
  \mathbb E_P[
    N_jE_{j,S\wedge(\tau_j+m)}\mid\mathcal F_{\tau_j}]
  \leq N_j.
  \label{eq:app-bounded-os}
\end{equation}
Now multiply the stopped variable by
\(\mathbf 1\{S\leq\tau_j+m\}\).  This can only decrease it, and the resulting
nonnegative variables converge pointwise to
\(N_j\overline E_{j,S}\): if \(S<\infty\), equality with
\(N_jE_{j,S}\) holds for all sufficiently large \(m\); if \(S=\infty\),
they are identically zero.  Conditional Fatou's lemma applied to
Eq.~\eqref{eq:app-bounded-os} proves
Eq.~\eqref{eq:app-selective-e}.  Taking expectations proves the final
statement.
\end{proof}

The product in Eq.~\eqref{eq:app-eprocess} is convenient rather than
essential.  For example, if under the null
\[
 \mathbb E_P[
   \exp\{\eta_{j,t}X_{j,t}-\eta_{j,t}^2v_{j,t}/2\}
   \mid\mathcal F_{t-1}]\leq1
\]
for predictable \(\eta_{j,t}\) and \(v_{j,t}\), the corresponding
exponential product is a global test supermartingale.  A convex mixture of
global e-processes is again globally valid.  Such constructions can replace
clipping when their conditional assumptions are economically defensible.

\section{The Irreversible Weighted Gate}
\label{app:gate}

Fix \(q\in(0,1)\).  Candidate \(j\) receives a nonnegative weight
\(\gamma_j\in\mathcal F_{\tau_j}\) such that
\begin{equation}
  \sum_{j=1}^{\infty}\gamma_j\leq1
  \quad\text{almost surely}.
  \label{eq:app-budget}
\end{equation}
The deterministic choice
\(\gamma_j=1/\{j(j+1)\}\) telescopes to one and accommodates an unbounded
proposal stream.  A proposal-time, data-dependent allocation is also
permitted, provided Eq.~\eqref{eq:app-budget} holds pathwise.  The convention
\(1/\gamma_j=\infty\) makes a zero-weight proposal ineligible for promotion.

At inspection \(t\), let \(\mathcal A_t\) be the finite,
\(\mathcal F_t\)-measurable set of eligible candidates satisfying
\begin{equation}
 D_{t-1}\subseteq\mathcal A_t\subseteq\{j:\tau_j<t\}.
 \label{eq:app-eligibility}
\end{equation}
Eligibility may include a prespecified post-proposal seasoning period, but an
unproposed index can never enter the gate.  Initialize
\(D_{-1}=\varnothing\), and suppose
inductively that \(D_{t-1}\) is finite and that each \(j\in D_{t-1}\) has a
frozen value
\[
  \widetilde E_j=E_{j,S_j},\qquad
  S_j=\inf\{s\geq\tau_j:j\in D_s\}.
\]
Form
\begin{equation}
 Z_{j,t}=
 \begin{cases}
   \widetilde E_j,&j\in D_{t-1},\\
   E_{j,t},&j\in\mathcal A_t\setminus D_{t-1}.
 \end{cases}
 \label{eq:app-gate-z}
\end{equation}
Set the mandatory individual evidence floor to \(h_q=q^{-1}\).  For
\(k\geq1\), let
\begin{equation}
 C_t(k)=\sum_{j\in\mathcal A_t}
   \mathbf 1\!\left\{
     Z_{j,t}\geq\max\!\left(h_q,\frac{1}{q\gamma_jk}\right)\right\}.
 \label{eq:app-floor-count}
\end{equation}
and define
\begin{align}
 k_t^\star
 &=\max\left(\{0\}\cup
    \{k\in\{1,\ldots,|\mathcal A_t|\}:C_t(k)\geq k\}\right),
 \label{eq:app-gate-k}\\
 D_t
 &=
 \begin{cases}
 \left\{j\in\mathcal A_t:
 Z_{j,t}\geq\max\!\left(h_q,(q\gamma_jk_t^\star)^{-1}\right)\right\},
       &k_t^\star>0,\\
 \varnothing,&k_t^\star=0.
 \end{cases}
 \label{eq:app-gate-d}
\end{align}
Every new member \(j\in D_t\setminus D_{t-1}\) receives \(S_j=t\) and
\(\widetilde E_j=E_{j,t}\).

\begin{proposition}[Well-definedness, irreversibility, and self-consistency]
\label{prop:app-gate}
The recursion in Eqs.~\eqref{eq:app-gate-z} and \eqref{eq:app-gate-d} is
adapted and finite.  For every \(t\),
\begin{equation}
  D_{t-1}\subseteq D_t,\qquad |D_t|=k_t^\star,
  \label{eq:app-nested-cardinality}
\end{equation}
and
\begin{equation}
  j\in D_t
  \quad\Longrightarrow\quad
  \widetilde E_j\geq h_q,\qquad
  \widetilde E_j\geq\frac{1}{q\gamma_j|D_t|}.
  \label{eq:app-self-consistency}
\end{equation}
Consequently, each \(S_j\) is a global stopping time.
\end{proposition}

\begin{proof}
All operations at inspection \(t\) involve a finite collection of
\(\mathcal F_t\)-measurable variables, so \(k_t^\star\) and \(D_t\) are
\(\mathcal F_t\)-measurable.  It remains to establish the invariants.
Proceed by induction.  They are immediate at initialization.  Write
\(d=|D_{t-1}|\).  If \(d>0\), the induction hypothesis gives, for every old
member,
\[
  \widetilde E_j\geq
  \max\!\left(h_q,\frac{1}{q\gamma_jd}\right).
\]
Hence \(C_t(d)\geq d\), so \(d\) is feasible in
Eq.~\eqref{eq:app-gate-k}.  If \(d=0\), feasibility of \(k=0\) is built into
the definition.  In either case \(k_t^\star\geq d\).  The map
\(k\mapsto\max\{h_q,(q\gamma_jk)^{-1}\}\) is nonincreasing, so every old
member exceeds its threshold at \(k_t^\star\).  Thus
\(D_{t-1}\subseteq D_t\).

If \(k_t^\star=0\), Eq.~\eqref{eq:app-gate-d} is empty and its cardinality
equals \(k_t^\star\).  Now suppose \(1\leq k_t^\star<|\mathcal A_t|\) and
\(C_t(k_t^\star)>k_t^\star\).  The threshold weakly decreases when \(k\)
increases, so
\[
 C_t(k_t^\star+1)\geq C_t(k_t^\star)\geq k_t^\star+1.
\]
This contradicts maximality.  If
\(k_t^\star=|\mathcal A_t|\), equality between the count and
\(k_t^\star\) is automatic.  Therefore the set in
Eq.~\eqref{eq:app-gate-d} has exactly \(k_t^\star\) members.  Its old members
retain their frozen values, and every new member is frozen at the value used
for admission.  This proves Eq.~\eqref{eq:app-self-consistency}.

Finally, nestedness and Eq.~\eqref{eq:app-eligibility} give
\(\{S_j\leq t\}=\{j\in D_t\}\in\mathcal F_t\).  Hence \(S_j\) is a global
stopping time.
\end{proof}

The floor is not required for joint FDR control, but it rules out a
weak-evidence passenger after other discoveries have lowered the step-up
boundary.  On the null branch, stopped e-validity and conditional Markov give
\(P(S_j<\infty\mid\mathcal F_{\tau_j})\leq q\).  This is a marginal statement
for candidate \(j\), not family-wise error control for the proposal stream.

\section{Simultaneous False-Discovery Control}
\label{app:supfdr}

The next result permits arbitrary dependence, distinct monitoring horizons,
adaptive proposal times, and infinitely many eventual proposals.  Its
multiple-testing step is the weighted
self-consistency principle behind e-BH and online e-BH
\cite{wang2022false,xu2024online,fischer2024online}; the proof must also
handle random strategy identities and global stopping of evolving evidence.

\begin{theorem}[Persistent SAFE-ALPHA certificate]
\label{thm:app-supfdr}
For each \(j\), suppose:
\begin{enumerate}
 \item \(\tau_j\) is a global stopping time and the candidate identity and
       null indicator are \(\mathcal F_{\tau_j}\)-measurable;
 \item \(E_{j,t}\) has the selective global e-validity property
       \eqref{eq:app-selective-e} at every global stopping time
       \(S\geq\tau_j\);
 \item \(\gamma_j\geq0\) is \(\mathcal F_{\tau_j}\)-measurable and the
       pathwise budget \eqref{eq:app-budget} holds; and
 \item \(D_t\) is a finite, adapted, nested discovery process satisfying
       Eq.~\eqref{eq:app-self-consistency}.
\end{enumerate}
Define
\[
 \operatorname{FDP}(D_t)
 =\frac{\sum_{j\in D_t}N_j}{|D_t|\vee1}.
\]
Then, under arbitrary serial and cross-candidate dependence,
\begin{equation}
 \mathbb E_P\!\left[
   \sup_{t\geq0}\operatorname{FDP}(D_t)\right]\leq q.
 \label{eq:app-supfdr}
\end{equation}
Consequently,
\begin{equation}
 \mathbb E_P[\operatorname{FDP}(D_T)]\leq q
 \label{eq:app-stopped-fdr}
\end{equation}
for every almost surely finite global stopping time \(T\).  The rule defining
\(T\) may use the complete research history, including market returns,
agent outputs, all evidence streams, portfolio performance, and prior gate
decisions.
\end{theorem}

\begin{proof}
Define the killed promotion value
\[
 \overline E_j
 =E_{j,S_j}\mathbf 1\{S_j<\infty\}.
\]
Proposition~\ref{prop:app-gate} makes \(S_j\) a global stopping time, and
selective e-validity gives
\begin{equation}
  \mathbb E_P[N_j\overline E_j\mid\mathcal F_{\tau_j}]
  \leq N_j.
  \label{eq:app-promotion-valid}
\end{equation}
Fix a sample path and a time \(t\).  If \(D_t=\varnothing\), its FDP is zero.
Otherwise, Eq.~\eqref{eq:app-self-consistency} implies, for every \(j\),
\[
 \frac{\mathbf 1\{j\in D_t\}}{|D_t|}
 \leq q\gamma_j\widetilde E_j\mathbf 1\{j\in D_t\}.
\]
Multiplication by \(N_j\), followed by summation, gives the pathwise bound
\begin{align}
 \operatorname{FDP}(D_t)
 &\leq q\sum_{j\in D_t}\gamma_jN_j\widetilde E_j
 \leq q\sum_{j=1}^{\infty}\gamma_jN_j\overline E_j.
 \label{eq:app-pathwise}
\end{align}
The last expression is independent of \(t\).  Therefore
\[
 \sup_{t\geq0}\operatorname{FDP}(D_t)
 \leq q\sum_{j=1}^{\infty}\gamma_jN_j\overline E_j.
\]
All summands are nonnegative.  Tonelli's theorem, proposal-time
measurability of \(\gamma_j\), and
Eq.~\eqref{eq:app-promotion-valid} now yield
\begin{align*}
 \mathbb E_P[\sup_t\operatorname{FDP}(D_t)]
 &\leq q\sum_{j=1}^{\infty}
    \mathbb E_P[\gamma_jN_j\overline E_j]\\
 &=q\sum_{j=1}^{\infty}
    \mathbb E_P\!\left[
      \gamma_j\,
      \mathbb E_P[N_j\overline E_j\mid\mathcal F_{\tau_j}]
    \right]\\
 &\leq q\sum_{j=1}^{\infty}\mathbb E_P[\gamma_jN_j]\\
 &\leq q\,\mathbb E_P\!\left[\sum_{j=1}^{\infty}\gamma_j\right]
 \leq q.
\end{align*}
This argument is unaffected by infinitely many eventual candidates:
finiteness at each inspection makes the gate well defined, and Tonelli
justifies the countable sum.  Finally, the pathwise inequality
\[
  \operatorname{FDP}(D_T)\leq\sup_t\operatorname{FDP}(D_t)
\]
proves Eq.~\eqref{eq:app-stopped-fdr}.
\end{proof}

The supremum appears \emph{inside} the expectation in
Eq.~\eqref{eq:app-supfdr}; this is stronger than fixed-time FDR control and
immediately implies control under optional stopping.  The argument depends
on freezing a single e-valid value at each global promotion time.  It does
not treat a candidate's raw historical maximum as an e-value.

\begin{proposition}[Containment of matched terminal e-BH]
\label{prop:app-terminal-containment}
At an inspection date \(t\), let \(R_t\) be the one-shot floor-augmented
weighted e-BH set obtained from the unfrozen endpoints \(E_{j,t}\) on
\(\mathcal A_t\), using the same \(q\), \(\gamma\), and \(h_q\) as the
persistent gate.  Then \(R_t\subseteq D_t\) pathwise.
\end{proposition}

\begin{proof}
Write \(k=|R_t|\) and \(d=|D_{t-1}|\).  Every terminal member passes the
floor and its size-\(k\) boundary.  By
Eq.~\eqref{eq:app-self-consistency}, every old member passes the floor and its
size-\(d\) boundary; every candidate outside \(D_{t-1}\) uses the same raw
endpoint in both procedures.  If \(k\leq d\), all terminal members also pass
at size \(d\), and \(k_t^\star\geq d\).  If \(k>d\), the union of the old and
terminal sets makes size \(k\) feasible, so \(k_t^\star\geq k\).  In either
case, the persistent threshold at \(k_t^\star\) is no larger than the boundary
already passed by each member of \(R_t\), proving the inclusion.
\end{proof}

\section{The Global-Filtration Condition and Its Boundaries}
\label{app:causal}

\subsection{A causal sufficient condition}

Equation~\eqref{eq:app-selective-null} is the exact condition used in the
proof.  The following representation gives it an operational interpretation.
Let \(C_t\in\mathcal F_{t-1}\) be a recorded market state and let \(Y_t\) be
the vector of contemporaneous innovations.  Suppose a transition kernel
\(K_t\) satisfies
\begin{equation}
  \mathcal L(Y_t\mid\mathcal F_{t-1})
  =K_t(\,\cdot\mid C_t).
  \label{eq:app-kernel}
\end{equation}
The committed strategy produces
\[
  X_{j,t}=g(A_j,C_t,Y_t).
\]
Because \(A_j\in\mathcal F_{t-1}\) whenever \(t>\tau_j\), the null condition
\begin{equation}
 N_j\int g(A_j,C_t,y)\,K_t(dy\mid C_t)\leq0
 \quad\text{on }\{\tau_j<t\}
 \label{eq:app-kernel-null}
\end{equation}
implies Eq.~\eqref{eq:app-selective-null} by conditional integration.  The
coordinates of \(Y_t\), and hence all strategy returns on the same date, may
be arbitrarily dependent.  The kernel conditions in
Eqs.~\eqref{eq:app-kernel} and \eqref{eq:app-kernel-null} do not assert a Markov
market.  They isolate the relevant causal restriction: no unrecorded
future innovation may reach the proposer, another evidence stream, or the
gate before the tested payoff is realized.  This is the global-filtration
condition emphasized for stopped e-BH by
\citet{wang2025stopped}.

The condition translates into auditable implementation rules:
\begin{enumerate}
 \item a signal using the close at \(t\) trades no earlier than the first
       implementable price after that close;
 \item proposal timestamps, data vintages, code hashes, random seeds, cost
       inputs, and execution lags are committed before evaluation;
 \item pre-proposal backtest returns never enter \(E_{j,t}\);
 \item a formula, universe, lag, or regime revision starts a new candidate;
 \item all evidence streams are aligned to the same physical market clock,
       rather than by observation number; and
 \item only the betting rule, not the tested payoff formula, may adapt
       predictably within an existing evidence process.
\end{enumerate}

\subsection{Same-bar leakage}
\label{app:same-bar}

Let \(X\) be a fair Rademacher payoff.  It satisfies the null
\(\mathbb E[X]=0\).  If a bettor illegally observes the same-bar payoff and
then sets
\[
  \lambda=\tfrac12\mathbf 1\{X=1\},
\]
the purported one-step evidence multiplier is \(3/2\) when \(X=1\) and \(1\)
when \(X=-1\).  Hence
\[
  \mathbb E[1+\lambda X]
  =\tfrac12(3/2)+\tfrac12(1)=\tfrac54>1.
\]
The loss of validity is entirely due to nonpredictability.  The example also
shows why a signal computed with a closing price cannot be credited with a
trade at that same closing price.

\subsection{Local validity does not imply global validity}
\label{app:local-global}

Let \(Y^1_1,Y^1_2,\ldots\) be independent Rademacher variables, and define a
second stream by \(Y^2_n=Y^1_{n+1}\).  In the filtration generated only by
the first stream,
\[
  M^1_n=\prod_{s=1}^n(1+Y^1_s/2)
\]
is a nonnegative martingale.  The global filtration containing both streams
reveals \(Y^1_{n+1}\) at time \(n\).  In that filtration,
\[
  S=1+\mathbf 1\{Y^2_1=1\}
    =1+\mathbf 1\{Y^1_2=1\}
\]
is a stopping time.  Nevertheless,
\begin{align*}
 \mathbb E[M^1_S]
 &=\tfrac12\mathbb E[M^1_2\mid Y^1_2=1]
   +\tfrac12\mathbb E[M^1_1\mid Y^1_2=-1]\\
 &=\tfrac12(3/2)+\tfrac12(1)=\tfrac54.
\end{align*}
Thus local optional-stopping validity is insufficient for a gate that
observes all streams.  This counterexample is the basic leakage mechanism
identified by \citet{wang2025stopped}.  SAFE-ALPHA requires every candidate's
e-process to be valid in the filtration in which the gate stops it.

\subsection{A random terminal family defeats ordinary e-BH}
\label{app:random-family}

Summable proposal weights are not a cosmetic replacement for the terminal
number of candidates.  Let all hypotheses be null and, independently for
\(j\geq1\), set
\[
 E_j=
 \begin{cases}
   j/q,&\text{with probability }q/j,\\
   0,&\text{otherwise}.
 \end{cases}
\]
Each \(E_j\) is an e-value because \(\mathbb E[E_j]=1\).  Stop proposing at
\[
  J=\inf\{j:E_j=j/q\}.
\]
Since \(\sum_jq/j=\infty\), independence implies \(J<\infty\) almost surely.
Ordinary unweighted e-BH applied to the random family
\(E_1,\ldots,E_J\) rejects \(H_J\): for one rejection its threshold is
\(J/q\), exactly \(E_J\).  It therefore has FDR one.  A deterministic
maximum \(M\) is safe with \(\gamma_j=1/M\) for all \(M\) slots, including
unused slots.  An unbounded adaptive stream requires a summable budget such
as Eq.~\eqref{eq:app-budget}.

\subsection{A raw running maximum is not an e-value}
\label{app:running-max}

Consider the one-step martingale \(E_0=1\) and
\[
  E_1=
  \begin{cases}
    2,&\text{with probability }1/2,\\
    0,&\text{with probability }1/2,
  \end{cases}
\]
held constant thereafter.  Although \(\mathbb E[E_1]=1\),
\[
  \mathbb E\!\left[\max_{s\leq1}E_s\right]
  =\tfrac12(2)+\tfrac12(1)=\tfrac32.
\]
The running maximum is therefore not an e-value.  Under dependence,
submitting raw e-process maxima to e-BH can inflate FDR; explicit
multiple-testing counterexamples are given by
\citet{tavyrikov2025carefree}.  SAFE-ALPHA does something different.  It
freezes \(E_{j,t}\) only when the \emph{current} value passes the joint
self-consistent gate.  That first-promotion time is global, and
Lemma~\ref{lem:app-selective-e} validates the single frozen value.  It never
recalls a prior sub-threshold peak.

If historical peaks are operationally indispensable, one must first apply a
valid adjuster.  A sufficient form is a nondecreasing function
\(A:[1,\infty]\to[0,\infty]\) satisfying
\[
  \int_1^\infty\frac{A(x)}{x^2}\,dx\leq1,
\]
and then use \(A(\sup_{s\leq t}E_{j,s})\).  The adjustment restores validity
but generally reduces power.

\section{Proposal-Rank Budget Allocation}
\label{app:rank-budget}

The spending weights determine the evidence boundary assigned to proposal
rank.  The following identity is an oracle calculation for the first useful
proposal; it is not a claim that the default telescoping schedule is generally
delay-optimal in a dynamic step-up process.

\begin{proposition}[Oracle one-discovery boundary]
\label{prop:app-oracle-budget}
Let \(J\) denote the rank of the first useful proposal and suppose
\(P(J=j)=\pi_j\).  Assume \(\sum_j\pi_j=\sum_j\gamma_j=1\) and
\(\gamma_j>0\) whenever \(\pi_j>0\).  Before another discovery has lowered
the step-up boundary,
\begin{equation}
 \mathbb E_\pi\!\left[\log\frac{1}{q\gamma_J}\right]
 =\log\frac1q+H(\pi)+\KL(\pi\,\|\,\gamma).
 \label{eq:app-rank-identity}
\end{equation}
The expectation is uniquely minimized by \(\gamma=\pi\).  If, in addition,
the log evidence grows deterministically as
\(\log E_{J,\tau_J+n}=rn\) for a common \(r>0\), then replacing the oracle
schedule by \(\gamma\) adds \(\KL(\pi\,\|\,\gamma)/r\) to expected
continuous crossing time.  On an integer inspection grid the equality holds
up to less than one inspection interval from rounding.
\end{proposition}

\begin{proof}
Direct expansion gives
\begin{align*}
 \mathbb E_\pi\!\left[\log\frac{1}{q\gamma_J}\right]
 &=\log\frac1q+\sum_j\pi_j\log\frac1{\gamma_j}\\
 &=\log\frac1q+H(\pi)
   +\sum_j\pi_j\log\frac{\pi_j}{\gamma_j}.
\end{align*}
The last sum is \(\KL(\pi\,\|\,\gamma)\), which is nonnegative and vanishes
exactly when \(\gamma=\pi\).  Under deterministic common log growth, crossing
time is the log boundary divided by \(r\), which proves the final statement.
\end{proof}

When proposal ranks and evidence growth rates are heterogeneous or dependent,
minimizing expected hitting time becomes a different design problem.  The
identity nevertheless quantifies the log-boundary regret of a misspecified
rank prior and explains why proposal-time economic information can be useful
without allowing post-proposal outcomes to determine the budget.

\section{Finite-Sample Power and Certification Delay}
\label{app:power}

The following bound makes the price of a late proposal explicit.  It is
finite-sample and allows arbitrary dependence compatible with the stated
conditional alternative.

\begin{proposition}[Certification by a fixed post-proposal horizon]
\label{prop:app-power}
Fix candidate \(j\).  Suppose that, for every \(s\geq1\),
\begin{equation}
 X_{j,\tau_j+s}\in[-1,1],
 \qquad
 \mathbb E_P[X_{j,\tau_j+s}\mid
              \mathcal F_{\tau_j+s-1}]\geq\mu
 \label{eq:app-alt}
\end{equation}
almost surely, where \(0<\mu\leq1\).  Use the constant predictable bet
\(\lambda=\mu/2\).  Set
\[
 g=\frac{\mu^2}{4},\qquad
 c_\mu=\log\frac{1+\mu/2}{1-\mu/2},
\]
and let \(d=|D_{\tau_j-1}|\), with \(D_{-1}=\varnothing\).  Conditional on
\(\mathcal F_{\tau_j}\), define
\begin{equation}
 L_j^{\rm floor}(d)
 =\log\frac1q+
  \left[\log\frac{1}{\gamma_j(d+1)}\right]_+.
 \label{eq:app-floor-delay-boundary}
\end{equation}
Suppose candidate \(j\) is eligible and the gate is inspected at
\(t=\tau_j+n\).  For any \(\delta\in(0,1)\), it is promoted no later than
that inspection with conditional probability at least \(1-\delta\) whenever
\begin{equation}
 ng-c_\mu\sqrt{\frac n2\log\frac1\delta}\geq L_j^{\rm floor}(d).
 \label{eq:app-delay-implicit}
\end{equation}
Equivalently, it suffices that
\begin{equation}
 n\geq
 \left[
 \frac{
 c_\mu\sqrt{\tfrac12\log(1/\delta)}
 +\sqrt{\tfrac12c_\mu^2\log(1/\delta)+4gL_j^{\rm floor}(d)}
 }{2g}
 \right]^2.
 \label{eq:app-delay-explicit}
\end{equation}
\end{proposition}

\begin{proof}
For \(|u|\leq1/2\), \(\log(1+u)\geq u-u^2\).  Put
\(Y_{j,t}=\log(1+\lambda X_{j,t})\).  Because
\(\lambda=\mu/2\leq1/2\), Eq.~\eqref{eq:app-alt} implies
\begin{align}
 \mathbb E_P[Y_{j,t}\mid\mathcal F_{t-1}]
 &\geq
 \lambda\mathbb E_P[X_{j,t}\mid\mathcal F_{t-1}]
 -\lambda^2\mathbb E_P[X_{j,t}^2\mid\mathcal F_{t-1}]
 \notag\\
 &\geq\lambda\mu-\lambda^2
 =\frac{\mu^2}{4}=g.
 \label{eq:app-log-drift}
\end{align}
Each log increment belongs to
\([\log(1-\lambda),\log(1+\lambda)]\), an interval of length \(c_\mu\).
Let
\[
 Z_s=Y_{j,\tau_j+s}
 -\mathbb E_P[Y_{j,\tau_j+s}\mid\mathcal F_{\tau_j+s-1}].
\]
These are martingale differences in the shifted filtration.  The conditional
Hoeffding lemma gives
\(\mathbb E_P[e^{-\eta Z_s}\mid\mathcal F_{\tau_j+s-1}]
\leq e^{\eta^2c_\mu^2/8}\).  Multiplication over \(s\), Markov's inequality,
and optimization over \(\eta>0\) yield
\[
 P\!\left(
   \left.\sum_{s=1}^nZ_s\leq-a\right|\mathcal F_{\tau_j}\right)
 \leq\exp\!\left(-\frac{2a^2}{nc_\mu^2}\right).
\]
Combining this inequality with Eq.~\eqref{eq:app-log-drift} and setting
\(a=c_\mu\sqrt{\tfrac n2\log(1/\delta)}\) gives
\begin{equation}
 P\!\left(
 \left.
 \log E_{j,\tau_j+n}
 <ng-c_\mu\sqrt{\frac n2\log\frac1\delta}
 \ \right|\mathcal F_{\tau_j}\right)\leq\delta.
 \label{eq:app-log-concentration}
\end{equation}

On the complementary event, Eq.~\eqref{eq:app-delay-implicit} yields
\[
E_{j,\tau_j+n}\geq
 \max\!\left\{\frac1q,\frac{1}{q\gamma_j(d+1)}\right\}.
\]
If \(j\) was promoted earlier, there is nothing to prove.  Otherwise let
\(m=|D_{\tau_j+n-1}|\).  Nestedness implies \(m\geq d\), and every one of
the \(m\) old discoveries remains above its threshold when the proposed set
size grows from \(m\) to \(m+1\).  Moreover,
\[
 E_{j,\tau_j+n}
 \geq\max\!\left\{\frac1q,\frac{1}{q\gamma_j(d+1)}\right\}
 \geq\max\!\left\{\frac1q,\frac{1}{q\gamma_j(m+1)}\right\}.
\]
Thus \(k=m+1\) is feasible in the maximal step-up rule, and \(j\) is
promoted at the current inspection.  This proves the probability statement.

To solve Eq.~\eqref{eq:app-delay-implicit}, put \(y=\sqrt n\) and
\(a=c_\mu\sqrt{\tfrac12\log(1/\delta)}\).  The required inequality is
\(gy^2-ay-L_j^{\rm floor}(d)\geq0\).  Its nonnegative root is
\[
 y=\frac{a+\sqrt{a^2+4gL_j^{\rm floor}(d)}}{2g},
\]
which gives Eq.~\eqref{eq:app-delay-explicit}.
\end{proof}

For \(\gamma_j=1/\{j(j+1)\}\),
\[
  L_j^{\rm floor}(d)=\log\frac1q+
  \left[\log\frac{j(j+1)}{d+1}\right]_+.
\]
Since \(c_\mu=O(\mu)\) and \(g=\mu^2/4\), the sufficient delay is
\begin{equation}
 O\!\left(
   \frac{\log j+\log(1/q)+\log(1/\delta)}{\mu^2}
 \right).
 \label{eq:app-delay-order}
\end{equation}
The penalty for proposal rank is logarithmic, while weak conditional drift
incurs the usual inverse-square delay.  A mixture over betting fractions or
a predictable plug-in bet avoids knowing \(\mu\) in advance, at a
logarithmic power cost, provided each component is a global e-process.

\section{From Certified Discoveries to Capital}
\label{app:capital}

The false-discovery certificate concerns a set of hypotheses.  It translates
directly into a capital statement only under an allocation rule that
preserves the relevant self-consistency.

\begin{corollary}[Equal weights]
\label{cor:app-equal-weight}
Suppose every member of \(D_t\) receives the same gross sleeve weight.  Set
the null-capital fraction to zero when \(D_t=\varnothing\).  Then
\begin{equation}
 \mathbb E_P\!\left[
   \sup_{t\geq0}
   \frac{\textnormal{gross capital in null sleeves at }t}
        {\textnormal{gross strategy capital at }t}
 \right]\leq q.
 \label{eq:app-equal-capital}
\end{equation}
\end{corollary}

\begin{proof}
Under equal gross sleeve weights, the displayed ratio is exactly
\(\sum_{j\in D_t}N_j/(|D_t|\vee1)=\operatorname{FDP}(D_t)\) on every sample
path.  Apply Theorem~\ref{thm:app-supfdr}.
\end{proof}

\begin{corollary}[E-budgeted allocation]
\label{cor:app-e-budget}
Let \(a_{j,t}\geq0\) be adapted gross-allocation shares satisfying
\begin{align}
 \sum_ja_{j,t}&\leq1,\notag\\
 a_{j,t}&=0 &&\text{for }j\notin D_t,\notag\\
 a_{j,t}&\leq q\gamma_j\widetilde E_j
       &&\text{for }j\in D_t.
 \label{eq:app-allocation-constraints}
\end{align}
Then
\begin{equation}
  \mathbb E_P\!\left[
    \sup_{t\geq0}\sum_{j=1}^{\infty}N_ja_{j,t}\right]\leq q.
  \label{eq:app-false-exposure}
\end{equation}
\end{corollary}

\begin{proof}
For every \(t\), pathwise,
\[
 \sum_jN_ja_{j,t}
 \leq q\sum_{j\in D_t}\gamma_jN_j\widetilde E_j
 \leq q\sum_{j=1}^{\infty}\gamma_jN_j\overline E_j.
\]
The right-hand side is independent of \(t\).  Take the supremum and repeat
the Tonelli and conditional-e-validity argument in the proof of
Theorem~\ref{thm:app-supfdr}.
\end{proof}

For example, a portfolio optimizer with gross capacity \(G\) may solve
\[
 \max_{w\in\mathcal W}
 \left\{
   \widehat\mu_t^\top w
   -\frac{\rho}{2}w^\top\widehat\Sigma_tw
   -c\lVert w-w_{t-1}\rVert_1
 \right\}
\]
subject to
\begin{equation}
 \operatorname{supp}(w)\subseteq D_t,\qquad
 \lVert w\rVert_1\leq G,\qquad
 \frac{|w_j|}{G}\leq q\gamma_j\widetilde E_j.
 \label{eq:app-optimizer-constraints}
\end{equation}
With \(a_{j,t}=|w_{j,t}|/G\), Corollary~\ref{cor:app-e-budget} controls false
gross exposure as a fraction of capacity.  It controls the fraction of
\emph{deployed} gross capital only when \(\lVert w_t\rVert_1=G\).  With
variable deployment, that interpretation requires defining
\(a_{j,t}=|w_{j,t}|/\lVert w_t\rVert_1\) and enforcing
Eq.~\eqref{eq:app-allocation-constraints} for those normalized shares.

Count-FDR alone does not justify arbitrary filtering after certification.
A downstream optimizer can place most capital in a false
member of an otherwise valid discovery set.  Nor does an e-value measure
effect size.  The certificate does not, by itself, control drawdown, CVaR,
turnover, realized Sharpe ratio, or estimation error in expected returns and
covariances.  Those quantities require a separate risk model and, where a
statistical guarantee is desired, confidence sequences or other
time-uniform estimators.  SAFE-ALPHA's role is narrower and auditable: it
limits false admission to the strategy set, and it limits false gross
exposure only when the allocation layer preserves the stated e-budget.
